\section{The Next Stage}
\label{sec:next-stage}

The next stage installs the current head as a filter, starts at the
following prime, and represents its own accepted sequence by a new
finite gap cycle. Throughout this section, a prime marks the next
stage's own version of each object defined in Section~2: $S'$ is the
next stage, $h'=\ell_1$ its head, $M'$ its tail primorial, $\ell'$ its
linear enumeration, and $G'$ its gap cycle.

\begin{itemize}
  \item the first old-stage successor is the next prime under the square
    bound;
  \item the semantic merged-gap prefix equals the next specification's
    gap prefix;
  \item a valid next-period boundary lets the next gap cycle reconstruct
    the next scan.
\end{itemize}

\subsection{Square-Bound Successor Primality}
\label{sec:square-bound-successor-primality}

The first old-stage successor is accepted by all filters smaller than
$h$. If that successor lies below $h^2$, then a composite successor
would have a prime divisor below $h$, contradicting acceptance.
Therefore the successor is prime under the square-bound precondition:

\begin{equation*}
\begin{aligned}
\ell_1&\lt h^2, \\
\ell_1\text{ composite}
&\Longrightarrow
\exists d\lt h,\ d\text{ prime and }d\mid \ell_1
  &&\text{[Smallest prime divisor]} \\
d\lt h
&\Longrightarrow d\in\overline{P}
  &&\text{[All smaller primes are filters]} \\
d\mid\ell_1
&\Longrightarrow \neg A_S(\ell_1)
  &&\text{[Filter contradiction]} \\
&\Longrightarrow \ell_1\text{ is prime}.
  &&\blacksquare\ \text{[Q.E.D.]}
\end{aligned}
\end{equation*}

{\raggedright
This property is verified in
\href{https://github.com/thiagomata/prime-numbers/blob/master/src/main/scala/v1/chapter6/sieve/seq/spec/properties/SpecSieveSeqHeadIsPrime.scala}{\texttt{SpecSieveSeqHeadIsPrime::assertApplyOneIsPrimeIfBelowHeadSq}}.
\par}

\subsection{The First Successor Is the Next Prime}
\label{sec:first-successor-is-next-prime}

Suppose the next prime after $h$ is $p^+$ and $p^+\lt h^2$. The next
prime passes every smaller-prime filter, so $\ell_1\le p^+$. Conversely,
if $\ell_1\lt h^2$ were composite, it would have a prime divisor at most
$\sqrt{\ell_1}\lt h$. That divisor belongs to $\overline{P}$,
contradicting acceptance. Thus $\ell_1$ is prime. No prime lies strictly
between $h$ and $p^+$, so $\ell_1=p^+$.

\begin{equation*}
\begin{aligned}
A_S(p^+) &\quad\text{[Distinct larger prime passes old filters]} \\
\ell_1 &\le p^+
  &&\text{[Least accepted successor]} \\
\ell_1 \lt h^2
  &\Longrightarrow \ell_1\text{ is prime}
  &&\text{[Small composite divisor]}
\end{aligned}
\end{equation*}

\begin{equation*}
h\lt\ell_1\le p^+,\quad \ell_1\text{ prime}
\ \Longrightarrow\ \ell_1=p^+
\quad\text{[No intervening prime]}\quad\blacksquare\ \text{[Q.E.D.]}
\end{equation*}

{\raggedright
This property is verified in
\href{https://github.com/thiagomata/prime-numbers/blob/master/src/main/scala/v1/chapter6/sieve/seq/spec/properties/SpecSieveSeqHeadIsPrime.scala}{\texttt{SpecSieveSeqHeadIsPrime::assertApplyOneEqualsNextPrime}}.
\par}

\subsection{Semantic Pipeline Agreement}
\label{sec:semantic-pipeline-agreement}

Let $T'=T(h-1)$. Under the stated stage-relationship invariants, the
semantic merge process starts at the first surviving old value and emits
$T'$ merged gaps. Section~5.3 supplies the copy-or-merge induction used
to prove that this list equals the first $T'$ gaps of the next linear
specification:

\begin{equation*}
\begin{aligned}
T'&=T(h-1), \\
\text{mergedGaps}(S,S',1,T')
&=\text{gapList}(S',0,T')
  \quad\text{[By copy-or-merge induction]}\quad\blacksquare\ \text{[Q.E.D.]}.
\end{aligned}
\end{equation*}

{\raggedright
This property is verified in
\href{https://github.com/thiagomata/prime-numbers/blob/master/src/main/scala/v1/chapter6/sieve/seq/spec/properties/SpecSieveSeqNextStageProperties.scala}{\texttt{SpecSieveSeqNextStageProperties::assertPipelineOutputMatchesNextGapList}}.
\par}

\subsection{Conditional Next-Cycle Reconstruction}
\label{sec:conditional-next-cycle-reconstruction}

If $T'$ is known to be the canonical period of the next stage, the
generic cycle-reconstruction theorem from Section~4.2 applies directly to
that stage:

\begin{equation*}
\begin{aligned}
\ell'_{T'} &= h'+M'
  &&\text{[Next-period boundary]} \\
G'&=\text{gapList}(S',0,T') \\
I_{G'}(k-1)&=\ell'_k
  &&\text{[Cycle reconstruction]}\quad\blacksquare\ \text{[Q.E.D.]}.
\end{aligned}
\end{equation*}

{\raggedright
This property is verified in
\href{https://github.com/thiagomata/prime-numbers/blob/master/src/main/scala/v1/chapter6/sieve/seq/spec/properties/SpecSieveSeqNextStageProperties.scala}{\texttt{SpecSieveSeqNextStageProperties::assertNextCycleReconstructsNextSpec}}.
\par}
